\documentclass[11pt]{article}
\usepackage[a4paper,margin=21mm]{geometry}
\usepackage[no-math]{fontspec}
\setmainfont{Latin Modern Roman}
\newfontfamily\CJKfont{Noto Serif CJK SC}
\XeTeXlinebreaklocale "zh"
\XeTeXlinebreakskip=0pt plus 1pt
\usepackage{amsmath,amssymb,amsthm,amscd,graphicx,color}
\usepackage{xurl}
\usepackage[unicode,hidelinks]{hyperref}
\allowdisplaybreaks
\setlength\emergencystretch{3em}
\numberwithin{equation}{section}


\numberwithin{equation}{section}

\newtheorem{Theorem}{Theorem}[section]
\newtheorem{Corollary}[Theorem]{Corollary}
\newtheorem{Lemma}[Theorem]{Lemma}
\newtheorem{Proposition}[Theorem]{Proposition}
 { \theoremstyle{definition}
\newtheorem{Definition}[Theorem]{Definition}
\newtheorem{Example}[Theorem]{Example}
\newtheorem{Remark}[Theorem]{Remark}
}

\newcommand{\e}{\mathrm{e}} % exponential e
\newcommand{\N}{\mathbb{N}}
\newcommand{\R}{\mathbb{R}}
\newcommand{\Z}{\mathbb{Z}}
\newcommand{\C}{\mathbb{C}}
\newcommand{\Q}{\mathbb{Q}}
\renewcommand{\P}{\mathbb{P}}
\newcommand{\E}{\mathbb{E}}
\newcommand{\dd}{\mathrm{d}} %integration d
\newcommand{\eps}{\varepsilon}
\newcommand{\la}{\langle}
\newcommand{\ra}{\rangle}

\newcommand{\be}{\begin{equation}}
\newcommand{\ee}{\end{equation}}

\newcommand{\vect}[1]{\boldsymbol{#1}}

\renewcommand{\Re}{\mathrm{Re}\,} %real part
\renewcommand{\Im}{\mathrm{Im}\,} %imaginary part
\DeclareMathOperator{\Int}{Int}
\DeclareMathOperator{\supp}{supp} % support
\DeclareMathOperator{\const}{const}

\newcommand{\set}[1]{{\left \{ #1 \right \}}}

\newcommand{\one}{\mathbbm{1}}
\newcommand{\1}{\mathbbm{1}}
\usepackage{mathrsfs}
\newfontfamily\nativebbone{DejaVu Math TeX Gyre}
\ExplSyntaxOn
\NewDocumentCommand{\mathbbm}{m}{\str_if_eq:nnTF{#1}{1}{\mathord{\text{\upshape\nativebbone\char"1D7D9}}}{\PackageError{nativecompat}{Unsupported~mathbbm~argument}{Only~verified~digit~one~is~accepted}}}
\ExplSyntaxOff
\renewcommand{\qedhere}{}

\begin{document}
\begin{center}\Large Unitary su(1,1) symmetries of consistent continuum particle Markov semigroups reversible for a finite-intensity Pascal law\end{center}
\noindent\textbf{Stable ID:} 986\quad\textbf{Author:} Simone Floreani; Sabine Jansen; Stefan Wagner\par
\noindent\textbf{Work:} Intertwinings for Continuum Particle Systems: an Algebraic Approach\par
\noindent\textbf{Source:} \url{https://sigma-journal.com/2024/046/}\par
\noindent\textbf{Version:} Official SIGMA 20 (2024) article 046, DOI 10.3842/SIGMA.2024.046; journal PDF and native TeX acquired 2026-09-30, exact bytes pinned by SHA256\par
\noindent\textbf{Rights:} CC BY 4.0. Authors retain copyright. \url{https://creativecommons.org/licenses/by/4.0/}. Reuse and typesetting adaptation; original language and mathematical content preserved.\par
\noindent\textbf{Difficulty (prior selection preserved):} {\CJKfont H2: The proof must control domains of unbounded creation and annihilation operators before passing semigroup commutation to their unitary exponentials. The bundled author argument establishes the Pascal Papangelou identity, adjointness, number conservation and essential self-adjointness through analytic-vector estimates, then proves domain invariance and spectral commutation. Formal Lie-algebra relations alone do not establish the conclusion.}\par
\medskip\noindent\textbf{Editorial boundary note (not author text):} Complete original Theorem 2.3 and its unbounded-domain/spectral-projection proof, including the full Pascal law and finite-configuration hypotheses, original operators and unitary definitions, Papangelou/Mecke argument, adjoint lemma, conservativity proof and entire analytic-vector essential-self-adjointness argument. The separate Meixner intertwining theorem and applications are excluded. The current algebra beyond the constant function is explicitly deferred by the author and is not a premise of this selected constant-function unitary symmetry. Only the verified double-struck digit one uses the named installed DejaVu Math TeX Gyre font; literal author mathbbm, one and 1 calls remain unchanged and other arguments are rejected. The single exact author qedhere call conflicts with the simultaneous source tag*\{qed\}; a disclosed QED-placement-only wrapper makes qedhere a no-op and retains the literal author proof-end tag and formula, following the original PDF.\par
\medskip\hrule\medskip\noindent\textbf{Author text begins below.}\par
\setcounter{section}{1}
% Original source lines 116--233
\section{Setting and main results} \label{sec:setting-results}

In the following, Cartesian products are always equipped with product $\sigma$-algebras. For finite interacting particle systems on lattices, the configuration space is $\N_0^\Lambda = \set{(n_x)_{x\in \Lambda}\colon n_x \in \N_0}$ with $\Lambda\subset\Z^d$. We replace $\N_0^\Lambda$ with a space of finite counting measures:
Let $(E,\mathcal E)$ be a Borel space, for example, $E= \R^d$ or more generally a Polish space with its Borel $\sigma$-algebra. Let~$\mathbf N_{<\infty}$ be the space of finite counting measures on $E$. Every element $\eta \in \mathbf N_{<\infty}$ is either zero or a finite sum $\eta= \delta_{x_1}+\cdots + \delta_{x_n}$ of Dirac measures, where $x_1, \ldots, x_n \in E$, $n \in \N$. The space $\mathbf N_{<\infty}$ is equipped with the $\sigma$-algebra $\mathcal N_{<\infty}$ generated by the counting variables $\eta \mapsto \eta(B)$, $B\in \mathcal E$.
For background on point processes and counting measures, we refer the reader to Last and Penrose~\cite{LastPenroseLecturesOnThePoissonProcess}.

\subsection{Symmetric inclusion process}

We are interested in Markov processes with state space $\mathbf N_{<\infty}$. As a guiding example we take the process with formal generator
\begin{equation} \label{eq:gsip}
	L f(\eta) = \iint \bigl( f(\eta-\delta_x+\delta_y)- f(\eta)\bigr) c(x,y) \eta(\dd x) \bigl( \alpha+\eta\bigr) (\dd y),
\end{equation}
where $\alpha$ is a finite measure on $E$ and $c\colon E\times E\to \R_+$ is a bounded measurable function with $c(x,y) = c(y,x)$ on $E^2$. In \cite{IntertwiningConsistent} we called this process \emph{generalized inclusion process}. When $E$ is a~finite set, we may identify finite counting measures $\eta$ with vectors $\vect n = (n_x)_{x\in \Lambda}\in \N_0^\Lambda$, integrals over $E$ turn into sums, and the generator turns into
\begin{equation} \label{eq:sip}
	Lf(\vect n) = \sum_{x,y\in E} \bigl( f(\vect n- \vect e_x+\vect e_y) - f(\vect n)\bigr) c(x,y) n_x (\alpha_y+n_y),
\end{equation}
where $\vect e_x(y) = \delta_{x,y}$ and $\alpha_y = \alpha(\{y\})$. This is the generator of the inhomogeneous \emph{symmetric inclusion process} (see, e.g., \cite[equation~(2.2)]{floreani2020orthogonal}), which first appeared in the homogeneous case as a dual process to a model for energy and momentum transport, see \cite{giardina-redig-vafayi2010} and references therein.

When $c(x,y)$ is the constant function with value $1$, the process corresponds to the Moran process from mathematical population genetics--classical Moran model \cite{TransitionMoran,Moran} when $E$ is finite, measure-valued Moran model (see \cite[Section 2.6]{DawsonMeasureValuedMarkovProcesses} or \cite[Section~5.4]{Etheridge00}) when $E$ is uncountable. The Moran model describes a population in which individuals have genetic types $x\in E$ that may mutate or evolve by sampling replacement (death of an individual followed by replacement with the offspring of another individual).

For the homogeneous symmetric inclusion process with finite state space / Moran model with finite type space, the correspondence was already noticed in \cite[Section~5]{carinci2015dualities}. For uncountable state spaces, we notice that the $n$-particle dynamics of the process generated by \eqref{eq:gsip} (with $c(x,y) \equiv 1$) coincides with the $n$-particle Moran model \cite[equation~(2.5.2)]{DawsonMeasureValuedMarkovProcesses} with so-called \emph{mutation operator} $A\varphi(x) = \int( \varphi(y) - \varphi(x)) \alpha(\mathrm d x)$. Our process for counting measures is similar to the measure-valued Moran model. The only difference is in bookkeeping: In mathematical population genetics, the population is often modeled not with $\eta \in \mathbf N_{<\infty}$ but rather with the empirical measure $\frac{\eta}{\eta(E)}$. This is helpful for scaling limits in which the population size goes to infinity, notably Fleming--Viot limits, which are important in population genetics (see \cite{shigaDiffusions} and references therein).

\subsection{Pascal point process} \label{sec:pascaldef}
The symmetric inclusion process with generator~\eqref{eq:gsip} has a family of reversible measures $\rho_{p,\alpha}$ indexed by $p\in (0,1)$ \cite[Theorem 5.2]{IntertwiningConsistent} where $\rho_{p,\alpha}$ is the law of the Pascal point process or negative binomial process (see, e.g., \cite[Section~2.7]{serfozo1990point}, \cite{kozubowski2008distributional} and \cite[Section~5.2]{IntertwiningConsistent}). We recall the definition here.
For $a\in \R$ and $n\in \N$, the rising factorial (also known as Pochhammer symbol)~is
\[
	(a)_0 = 1,\qquad (a)_n = a(a+1)\cdots (a+n-1).
\]

\begin{Definition}
 Let $p\in (0,1)$ and $\alpha$ be a finite measure on $E$. $\rho_{p,\alpha}$ is the uniquely defined probability measure on $\mathbf N_{<\infty}$ such that:
 \begin{itemize}\itemsep=0pt
 	\item For all $B\in \mathcal E$, and $n\in \N_0$
 	\begin{equation} \label{eq:negbin}
 		\rho_{p,\alpha}(\{\eta\colon \eta(B) =n\}) = (1-p)^{\alpha(B)} \frac{p^n}{n!} (\alpha(B))_{n}.
 	\end{equation}
 	\item For all $n \geq 2$ and all disjoint $B_1,\ldots, B_n\in \mathcal E$, the variables $\eta\mapsto \eta(B_i)$, $i=1,\ldots, n$, are independent.
 \end{itemize}
\end{Definition}

We work in the Hilbert space $\mathcal H = L^2(\mathbf N_{<\infty}, \mathcal N_{<\infty}, \rho_{p,\alpha})$ of complex-valued square-integrable functions. The scalar product is
\[
	\langle f,g \rangle = \int \overline{f} g\, \dd\rho_{p,\alpha}.
\]

\subsection[Representation of the su(1,1)]{Representation of the $\boldsymbol{\mathfrak{su}(1,1)}$ current algebra} \label{sec:su11}

Let $\mathcal D\subset \mathcal H$ be the set of bounded measurable functions $F$ for which there exists a cutoff $ m_F \in \N$ such that $F$ is supported in $\{\eta\colon \eta(E)\leq m_F\}$. We define operators $k^\pm(\varphi),k^0(\varphi)\colon\mathcal D\to \mathcal D$ indexed by bounded measurable functions $\varphi\colon E\to \C$ as follows:
\begin{align*}
	&k^+(\varphi) F(\eta) = \frac{1}{\sqrt p} \int \varphi(x) F(\eta-\delta_x) \eta(\dd x),\\
 &k^-(\varphi) F(\eta) = \sqrt{p} \int \overline{\varphi(x)} F(\eta+\delta_x) \bigl(\alpha+\eta)\bigr(\dd x),\\
	&k^0(\varphi) F(\eta) = F(\eta) \biggl( \int \varphi\,\dd \eta+ \frac12 \int \varphi \,\dd \alpha \biggr).
\end{align*}
for $\eta \in \mathbf{N}_{<\infty}$.
The raising operator $k^+(\varphi)$ and the lowering operator $k^-(\varphi)$ map functions supported in $\{\eta\colon \eta(E) = n\}$ to functions supported in $\{\eta\colon \eta(E) =n+1\}$ and $\{\eta\colon \eta(E) =n-1\}$ respectively. The indicator $\1_{\{0\}}$ that there is no particle at all (\emph{vacuum}) is annihilated by all lowering operators: $k^-(\varphi) \1_{\{0\}}=0$.

Our operators are a representation of the $\mathfrak{su}(1,1)$ current algebra, i.e., they represent the matrix-valued maps $x\mapsto \varphi(x) k^+$, $x\mapsto \overline{\varphi(x)} k^-$, $x\mapsto \theta(x) k^0$ with
\begin{align*}
 	k^+ = \begin{pmatrix} 0 & \mathrm i \\ 0 & 0 \end{pmatrix},\qquad
	k^- = \begin{pmatrix} 0 & 0 \\ \mathrm i & 0 \end{pmatrix},\qquad
	k^0 = \begin{pmatrix} 1/2 & \hphantom{-}0 \\ 0 & -1/2 \end{pmatrix},
\end{align*}
see \cite[Section~2.2]{algebraic_Floreani}.
The following commutation relations can be checked by an explicit computation:
\begin{gather*}
	\bigl[k^-(\varphi), k^+(\theta)\bigr] = 2 k^0(\overline{\varphi}\theta),\qquad
\bigl[k^0(\theta),k^+(\varphi)\bigr] = k^+(\varphi\theta),\qquad [k^0(\theta),k^-(\varphi)] = - k^-(\overline{\varphi}\theta),
\end{gather*}
where $[T, S] := TS - ST$.
Indeed, it can be readily verified that
\begin{gather*}
 k^-(\varphi) k^+(\theta) F(\eta) = \iint \overline{\varphi(y)} \theta(x) F(\eta - \delta_x + \delta_y) \eta(\dd x) (\alpha + \eta)(\dd y) \\
 \hphantom{k^-(\varphi) k^+(\theta) F(\eta) = \iint}{}
 + F(\eta) \int \overline{\varphi(y)} \theta(y) (\alpha + \eta)(\dd y), \\
 k^+(\theta) k^-(\varphi) F(\eta) = \iint \overline{\varphi(y)} \theta(x) F(\eta - \delta_x + \delta_y) \eta(\dd x) (\alpha + \eta)(\dd y) \\
 \hphantom{k^+(\theta) k^-(\varphi) F(\eta) =\iint}{}
 + F(\eta) \int \overline{\varphi(x)} \theta(x) \eta(\dd y).
\end{gather*}
Subtracting these equations yields the commutation relation
\[
\bigl[k^-(\varphi), k^+(\theta)\bigr] F(\eta) = F(\eta) \int \overline{\varphi(x)} \theta(x) (\alpha + 2\eta)(\dd x) = 2k^0(\overline{\varphi}\theta) F(\eta).
\]
 The other relations follow similarly.

The factor $\sqrt p$ included in the definitions of $k^\pm(\varphi)$ is irrelevant for the commutation relations but it matters for the adjointness relation $\la f, k^+(\varphi) g\ra = \la k^-(\varphi) f,g\ra$ proven in Lemma~\ref{lem:adjoints}.

We leave a thorough analysis of the current algebra generated by the operators $k^\pm(\varphi)$, $k^0(\varphi)$ to another article \cite{algebraic_Floreani} and focus on the operators associated with the constant function $\varphi = \1$, equal to $1$ on all of $E$. In Lemma~\ref{lem:closure} below, we check that for every $\xi \in \C$, the operator $\frac{1}{\mathrm i} \bigl( \xi k^+(\1) - \overline{\xi} k^-(\1))$ with domain $\mathcal D$ is closable with self-adjoint closure $A$. The operator $\exp( \xi k^+(\1) - \overline{\xi} k^-(\1))$ is, by definition, the unitary $\exp(\mathrm i A)$.
For $\theta \in \R$, the unitary $\exp\bigl( 2\mathrm i \theta k^0(\1)\bigr)$ is defined as a multiplication operator, it multiplies a function $f(\eta)$ with $\exp(\mathrm i \theta (\alpha (E) + 2 \eta(E)))$.
In this way we obtain a family of unitary operators
\begin{equation}\label{eq:unitary}
	{\rm U}(\xi,\theta) = \exp\bigl( \xi k^+(\1) - \overline{\xi} k^-(\1)\bigr) \exp\bigl(2\mathrm i\theta k^0(\1)\bigr),\qquad \xi \in \C,\theta\in \R.
\end{equation}
For the definition of the unitary~\eqref{eq:unitary}, it is essential that $\alpha$ has finite total mass and that each configuration $\eta$ has finitely many particles.

\subsection{Symmetries of consistent particle processes}

Let $(P_t)_{t\geq 0}$ be the semigroup of a Markov process with state space $\mathbf N_{<\infty}$. For a measurable $f\colon \mathbf N_{<\infty} \to \R_+$, let $\mathcal A f\colon \mathbf N_{<\infty} \to \R_+$ be the function given by
\[
	\mathcal A f(\eta) = \int f(\eta-\delta_x)\eta(\dd x).
\]

\begin{Definition}
 The semigroup $(P_t)_{t\geq 0}$ is \emph{consistent} if $P_t \mathcal A f = \mathcal A P_t f$ for all measurable $f\colon \mathbf N_{<\infty}\to \R_+$ and all $t\geq 0$.
\end{Definition}

Notice that, up to questions of domains, $\mathcal A = \sqrt p\, k^+(\1)$. In Proposition~\ref{proposition: consistency implies conservative} below, we show that if a process is consistent then it is \emph{conservative}, i.e., the total number of particles is constant in time. In view of that, we could normalize $\mathcal A$ by the total mass $\eta(E)$ (which is constant in time) and read consistency as the property that the action of removing a particle uniformly at random commutes with the dynamics \cite{carinci2021consistent,IntertwiningConsistent}.

The following theorem generalizes item 1\,(i) in Theorem~3.1 in Carinci et al.~\cite{carinci2019orthogonal}.

\begin{Theorem} \label{thm:symmetry}
Let $(P_t)_{t\geq 0}$ be the semigroup of a Markov process with state space $\mathbf N_{<\infty}$. Let $p\in (0,1)$ and $\alpha$ be a finite measure on $E$. Assume that the process is consistent, and admits the Pascal law $\rho_{p,\alpha}$ as a reversible measure. Then $P_t$ commutes with all unitaries ${\rm U}(\xi,\theta)$ from~\eqref{eq:unitary}:
	\[
		P_t {\rm U}(\xi,\theta) f(\eta) = {\rm U}(\xi,\theta) P_t f (\eta)
	\]
	for all $t\geq 0$, $f\in L^2(\mathbf N_{<\infty},\mathcal N_{<\infty},\rho_{p,\alpha})$ and $\rho_{p,\alpha}$-almost all $\eta\in \mathbf N_{<\infty}$.
\end{Theorem}

The generalized symmetric inclusion process is a consistent Markov process and admits the Pascal law $\rho_{p,\alpha}$ as a reversible measure (as proved in \cite[Theorem~5.2.]{IntertwiningConsistent}) and, thus, Theorem~\ref{thm:symmetry} applies.

\setcounter{section}{2}
% Original source lines 313--472
\section{Papangelou kernel for the Pascal point process} \label{sec:pascal}


If $X$ is a negative binomial variable with parameters $p\in (0,1)$ and $a>0$, then
\[
 (n+1)\P( X = n+1) = (1-p)^a\, \frac{(a)_{n+1}}{n!} p^{n+1} = p (a+n)\P( X = n),
\]
for all $n\in \N_0$, accordingly
\[
	\E[ X f(X)] = \E[ p(a+X) f(X+1)]
\]
for all $f\colon \N_0\to \R_+$. The following proposition gives the analogous property for Pascal point processes. For future reference we state and prove the proposition for $\sigma$-finite measures $\alpha$, allowing for configurations with infinitely many particles. Thus, let $\mathbf N$ be the space of $s$-finite counting measures on $E$. As $E$ is a Borel space, $\mathbf N$ consists of the measures $\eta$ on $E$ that are finite or countable sums of Dirac measures. The space $\mathbf N$ is equipped with the $\sigma$-algebra $\mathcal N$ generated by the counting variables $B\mapsto \eta(B)$, $B\in \mathcal E$.

A \emph{Pascal point process} with parameters $p$ and $\alpha$ is a random variable $\eta\colon(\Omega,\mathcal F, \P)\to (\mathbf N,\mathcal N)$ that satisfies the properties listed in the definition of $\rho_{p,\alpha}$: Counting variables for disjoint regions~$B_i$ are independent, and each counting variable $\eta(B)$ has a negative binomial distribution with parameters $p$ and $\alpha(B)$. For sets with infinite mass $\alpha(B) = \infty$, the counting variable $\eta(B)$ is almost surely infinite. Pascal point processes have Laplace functional
\begin{gather*}
	\E\bigl[ \e^{-\int f\dd \eta}\bigr] = \exp \biggl(\! - \int \log \biggl(\frac{1- p\,\e^{-f(x)}}{1- p}\biggr) \alpha (\dd x)\biggr) = 	\exp\Biggl(\! - \sum_{j=1}^\infty \int \bigl(1- \e^{- j f(x)}\bigr) \frac{p^j}{j} \alpha(\dd x) \Biggr)
\end{gather*}
for measurable $f \colon E \to \R_+$.

\begin{Proposition}%\label{prop:papangelou}
	Fix $p\in (0,1)$ and a $\sigma$-finite measure $\alpha$ on $E$. Let $\eta$ be a Pascal point process with parameters $p$ and $\alpha$. Then
	\[
		\E\biggl[ \int F(x,\eta)\eta(\dd x)\biggr] = \E\biggl[\int F(x,\eta+\delta_x) p (\alpha+\eta)(\dd x)\biggr]
	\]
	for all measurable $F\colon E\times \mathbf N \to \R_+$.
\end{Proposition}

The proposition says that the kernel $E\times \mathcal N\to \R_+$, $(x,B)\mapsto \kappa(x,B) = p(\alpha(B) + \eta(B))$ is a~\emph{Papangelou kernel} for the Pascal point process. For a general discussion on Papangelou kernels, we refer to \cite{MatthesWarmuthMecke} or \cite[Section 1.2.2]{Rafler2009}.

\begin{proof}
	Pascal point processes are compound Poisson; we deduce the proposition from the Mecke formula for Poisson point processes. Let $\Pi = \sum_i \delta_{(X_i,N_i)}$ be a Poisson point process on $E\times \N_0$ with intensity measure $\beta = \alpha\otimes \bigl(\sum_{n=1}^\infty \frac{p^n}{n} \delta_n\bigr)$. Then $\xi = \sum_i N_i \delta_{X_i}$ is a Pascal point process with law $\rho_{p,\alpha}$. The Mecke equation for $\Pi$ \cite[Theorem 4.1]{LastPenroseLecturesOnThePoissonProcess} and measurable $G\colon (E\times \N_0)\times \mathbf N(E\times \N_0) \to \R_+$ gives
\[
	\E\biggl[ \sum_i G( (X_i, N_i), \Pi )\biggr] = \sum_{n=1}^\infty \frac{p^n}{n} \int \E\bigl[ G\bigl( (x,n), \delta_{(x,n)} + \Pi\bigr) \bigr] \alpha (\dd x).
\]
Let $\varphi\colon E\times \mathbf N(E)\to \R_+$ be a measurable function. We apply the Mecke equation for $\Pi$ to $G\bigl((x,n),\sum_i \delta_{(x_i,n_i)}\bigr):= n \varphi\bigl(x, \sum_i n_i \delta_{x_i}\bigr)$ and get
\begin{equation}
 \label{equation: mecke1}
	\E\biggl[\sum_i N_i \varphi( X_i, \xi)\biggr] = \sum_{n=1}^\infty p^n \int \E[ \varphi(x, \xi+ n \delta_x)] \alpha(\dd x).
\end{equation}
Using \eqref{equation: mecke1} for $\varphi(x,\mu) = F(x,\mu + \delta_x)$, we obtain
\begin{align}
 p\, \E \biggl[ \sum_i N_i F\bigl( X_i, \xi +\delta_{X_i}\bigr)\biggr] &= p\sum_{n=1}^\infty p^n \int \E [ F(x, \xi + (n+1) \delta_{x}) ]\alpha(\dd x) \notag \\
 &= \sum_{n=2}^\infty p^n \int \E [ F(x, \xi + n \delta_{x}) ] \alpha(\dd x). \label{eq:mecke2}
\end{align}
We apply~\eqref{equation: mecke1} to $\varphi(x,\mu) = F(x,\mu)$ and eliminate the sum over $n\geq 2$ using~\eqref{eq:mecke2}. This gives
\begin{equation*}
	\E\biggl[\sum_i N_i \varphi( X_i, \xi)\biggr] =
		p \int \E[ F(x, \xi+ \delta_x)] \alpha(\dd x) + p\, \E \biggl[ \sum_i N_i F\bigl( X_i, \xi +\delta_{X_i}\bigr)\biggr].
\end{equation*}
Equivalently,
\[
	\E\biggl[ \int F(x,\xi) \xi(\dd x)\biggr] = p\, \E\biggl[ \int F(x, \xi+\delta_x) (\alpha + \xi)(\dd x)\biggr]
\]
and the proof is complete as $\xi$ and $\eta$ are equal in distribution.
\end{proof}

Turning back to the law $\rho_{p,\alpha}$ on $\mathbf N_{<\infty}$ for finite measures $\alpha$, we obtain
\begin{equation} \label{eq:papangelou}
	\int \biggl( \int F(x,\eta)\eta(\dd x)\biggr) \rho_{p,\alpha}(\dd x)
	= \int\biggl( \int F(x,\eta+\delta_x) p (\alpha+\eta)(\dd x)\biggr) \rho_{p,\alpha}(\dd x)
\end{equation}
for all measurable $F\colon E\times \mathbf N_{<\infty}\to \R_+$, hence also all bounded measurable functions $F$ (here we use that the expected total number of particles is finite when $\alpha$ has finite total mass).

\begin{Lemma} \label{lem:adjoints}
	For all bounded measurable $\varphi\colon E\to \C$ and all $f,g\in \mathcal D$,
	\[
		\bigl\langle f, k^+(\varphi) g\bigr\rangle = \langle k^-(\varphi) f,g \rangle.
	\]
\end{Lemma}

\begin{proof}	
	Let us denote integration with respect to $\rho_{p,\alpha}$ as $\E[f(\eta)] =\int f\dd \rho_{p,\alpha}$. We apply equation~\eqref{eq:papangelou} to $F(x,\eta) = \varphi(x) \overline{f(\eta)} g(\eta - \delta_x)$ and obtain
	\begin{align*}
		\bigl\langle f, k^+(\varphi) g \bigr\rangle
			& = \frac 1{\sqrt p} \E\biggl[\int \varphi(x) \overline{f(\eta)} g(\eta - \delta_x) \eta(\dd x)\biggr]\\
			& = \sqrt p\, \E\biggl[\int \varphi(x) \overline{f(\eta + \delta_x)} g(\eta) (\alpha+ \eta)(\dd x)\biggr]\\
			& = \langle k^-(\varphi) f,g \rangle. \qedhere
\tag*{\qed}
\end{align*}
\renewcommand{\qed}{}
\end{proof}

\section{Symmetries. Proof of Theorem~\ref{thm:symmetry}} \label{sec:symmetry}

Before we turn to the proof of Theorem~\ref{thm:symmetry}, we shows that consistency implies conservativity, which is of interest in its own.

\begin{Proposition}
 \label{proposition: consistency implies conservative}
	Every consistent Markov semigroup is conservative.
\end{Proposition}

\begin{proof}
	Let $(P_t)_{t\geq 0}$ be Markov semigroup on $(\mathbf N_{<\infty}, \mathcal N_{<\infty})$ that is consistent, i.e., $P_t \mathcal A f = \mathcal A P_t f$ for all $t\geq 0$ and all measurable $f\colon \mathbf N_{<\infty} \to \R_+$. Then, for all $k\in \N$,
	\begin{equation} \label{eq:ptak}
		P_t \mathcal A^k \one = \mathcal A^k P_t \one = \mathcal A^k \one.
	\end{equation}
	A straightforward computation shows that $\mathcal A^k\one$ is a descending factorial power of the total number of particles,
	\begin{equation} \label{eq:ak-factorial}
		\bigl( \mathcal A^k \one \bigr)(\eta) = \eta(E)( \eta(E)-1) \cdots ( \eta(E) - k +1),
	\end{equation}
	for all $\eta \in \mathbf N_{<\infty}$. Let $(\eta_t)_{t\geq 0}$ be a Markov process with semigroup $(P_t)_{t\geq 0}$ and deterministic initial condition $\eta_0=\mu \in \mathbf N_{<\infty}$, defined on some probability space $\bigl(\Omega,\mathscr F, \mathbb P_\mu\bigr)$. Equations~\eqref{eq:ptak} and~\eqref{eq:ak-factorial} show that under $\mathbb P_\mu$, at all times $t\geq 0$, the total number of particles $\eta_t(E)$ has the same factorial moments as $\eta_0(E)$. The factorial moments of $\eta_0(E) = \mu(E)$ (hence $\eta_t(E)$) vanish for $k\geq \mu(E)+1$, therefore the moment problem is uniquely solvable and we conclude $\eta_t(E) = \eta_0(E) = \mu(E)$, $\P_\mu$-almost surely, and the process is conservative.
\end{proof}

Remember that the set $\mathcal D\subset \mathcal H$ consists of the bounded measurable functions $f$ supported in $\{\eta\colon \eta(E)\leq n_f\}$ for some $n_f<\infty$.

\begin{Lemma} \label{lem:closure}
	For every $\xi \in \C$, the operator $\frac{1}{\mathrm i} \bigl( \xi k^+(\1) - \overline{\xi} k^-(\1)\bigr)$ with domain $\mathcal D$ is closable and its closure is self-adjoint.
\end{Lemma}

Put differently, the operator is \emph{essentially self-adjoint} on $\mathcal D$ and the latter is a \emph{core} for the self-adjoint closure \cite[Section VIII.2]{reedSimonI}.

\begin{proof}
	The operator $A = \frac{1}{\mathrm i} \bigl( \xi k^+(\1)- \overline{\xi} k^-(\1)\bigr)$ is similar to the Segal field operators for free quantum fields, we adapt the proof of essential self-adjointness for the latter from Reed and Simon \cite[Theorem X.41\,(a)]{reedSimonII}.
	The domain $\mathcal D$ is dense in $\mathcal H$ and satisfies $\la f, A g\ra = \la Af, g\ra$ for all $f,g\in \mathcal D$ by Lemma~\ref{lem:adjoints}. Thus, $A$ with domain $\mathcal D(A) = \mathcal D$ is symmetric. Symmetric operators are always closable. 	
	To show that the closure is self-adjoint, we check that every vector $f\in \mathcal D$ is \emph{analytic} for $A$ and conclude with Nelson's analytic vector theorem \cite[Theorem~X.39]{reedSimonII}.
 As a reminder, Nelson's analytic vector theorem states that a symmetric operator on a Hilbert space whose domain contains a total set of analytic vectors is essential self-adjoint.
 A vector~$f$ is analytic for $A$ if there exists $\eps>0$ such that $\sum_{n=0}^\infty n!^{-1} \eps^n \| A^n \varphi\|<\infty$.
	Let $f\in \mathcal D$ and $m=m_f\in \N$ be such that $|f(\eta)|\leq \|f\|_\infty \1_{\{\eta(E)\leq m_f\}}$ on $\mathbf N_{<\infty}$. Then
	\begin{align*}
		&\bigl| k^+(\1) f(\eta)\bigr| \leq \frac{1}{\sqrt p} \|f\|_\infty (m_f+1) \1_{\{\eta(E) \leq m_f+1\}},\\
		&| k^-(\1) f(\eta)| \leq \sqrt p\, \|f\|_\infty \bigl(\alpha(E)+m_f-1\bigr) \1_{\{\eta(E) \leq m_f-1\}}
	\end{align*} 	
	for all $\eta \in \mathbf N_{<\infty}$. Set $\beta:= \alpha(E)$, we obtain a common upper bound to $k^\pm(\1) f$:
	\[
		\bigl| k^\pm(\1) f(\eta)\bigr| \leq \frac{1}{\sqrt p} \|f\|_\infty (\beta +m_f+1) \1_{\{\eta(E) \leq m_f + 1\}}.
	\]
	We expand the power $\bigl(k^+(\1) - k^-(\1)\bigr)^n f$ and bound each term in the expansion by repeated use of the previous inequality. This gives
	\[
		\bigl| \bigl(k^+(\1) - k^-(\1)\bigr)^n f(\eta)\bigr| \leq \biggl(\frac{ 2 \| f\|_\infty}{\sqrt p}\biggr)^{n} ( \beta + m_f + 1)_n \1_{\{\eta(E)\leq m_f+n\}}.
	\] 	
	The indicator on the right has norm $\leq 1$. Set $s = 2|\xi|\, \|f\|_\infty/\sqrt p$, we obtain
	\[
		\sum_{n=0}^\infty \frac{\eps^n}{n!} \bigl\| \bigl(\xi k^+(\1) -\overline{\xi} k^-(\1)\bigr)^n f \bigr\|
			 \leq \sum_{n=0}^\infty \frac{(s\eps)^n}{n!} \bigl(\beta + m_f + 1)_n,
	\] 	
	which is finite for $\eps< 1/s$. Thus, every vector $f\in \mathcal D$ is analytic and $A$ is essentially self-adjoint on $\mathcal D$. 		
\end{proof}

\begin{proof}[Proof of Theorem~\ref{thm:symmetry}]
	Let $t \geq 0$. As $\rho_{p,\alpha}$ is reversible, $P_t$ is a well-defined, self-adjoint and bounded operator on $L^2(\mathbf N_{<\infty}, \mathcal N_{<\infty}, \rho_{p, \alpha})$. The operator $\exp\bigl(2 \mathrm i \theta k^0(\1)\bigr)$ acts as multiplication with $\exp( \mathrm i \theta (\alpha(E) + 2\eta(E)))$, it commutes with $P_t$ because $(P_t)_{t \geq 0}$ preserves the total number of particles (see Proposition~\ref{proposition: consistency implies conservative} below).
	
		For $\exp\bigl( \xi k^+(\1) - \overline{\xi} k^-(\1)\bigr)$, the short informal reasoning is simple: $P_t$ is self-adjoint and commutes with $k^+(\1)$, therefore it also commutes with the adjoint $k^-(\1)$ and also with the difference $\xi k^+(\1) - \overline{\xi} k^-(\1)$ and the exponential $\exp\bigl( \xi k^+(\1) - \overline{\xi} k^-(\1)\bigr)$. The precise reasoning is slightly longer because the operators are unbounded and we have to consider domains carefully. Let $f$, $g$ be non-negative functions in $\mathcal D$. Arguing as in the proof of Lemma~\ref{lem:adjoints}, we obtain $\langle k^-(\1) f, P_t g\rangle = \langle f, \mathcal A P_t g \rangle$. Therefore,
	\[
		\langle k^-(\1) f, P_t g\rangle = \langle f, \mathcal A P_t g \rangle = \langle f, P_t \mathcal A g\rangle = \bigl\langle P_t f, k^+(\1) g\bigr\rangle
	\]
	for all non-negative $f,g\in \mathcal D$. 	Taking complex conjugates and switching $f$ and $g$, we get \smash{$\bigl\langle k^+(\1) f, P_t g \bigr\rangle = \langle P_t f, k^-(\1) g\rangle $} and then, by linear combination,
	\[
		\biggl \langle \frac 1{\mathrm i} \bigl(\xi k^+(\1) - \overline{\xi} k^-(\1)\bigr) f, P_t g\biggr \rangle =
		\biggl \langle P_t f, \frac 1{\mathrm i} \bigl(\xi k^+(\1) - \overline{\xi} k^-(\1)\bigr) g\biggr \rangle
	\]
	for all $f,g\in \mathcal D$. By Lemma~\ref{lem:closure}, the operator $- \mathrm {i}\bigl(\xi k^+(\1) - \overline{\xi} k^-(\1) \bigr)$ with domain $\mathcal D$ is closable and its closure $A$ is self-adjoint. 	
	Taking limits in the above equation, we obtain
	\[
		\langle A f, P_t g \rangle = \langle P_t f, A g\rangle
	\]
	for all $f$, $g$ in the domain of $A$. In particular, for fixed $g$ in the domain of $\mathcal A$, and all $f\in \mathcal D(A)$, we have $\langle A f, P_t g\rangle = \langle f, P_t A g \rangle$, hence $P_tg$ is in the domain of $A^*$ and $A^* P_t g = P_t A g$.
	
	As $A$ is self-adjoint, we conclude that $P_t$ leaves the domain of $A$ invariant and $P_t A = A P_t$ on the domain of $A$.
It follows that $P_t$ commutes with all spectral projections of $A$ \cite[Proposition~5.26]{schmudgen2012unbounded} and then, by spectral calculus, $P_t \exp(\mathrm i A) = \exp(\mathrm i A) P_t$. 	
\end{proof}
\begin{thebibliography}{99}
\bibitem[7]{carinci2019orthogonal}
Carinci G., Franceschini C., Giardin\`a C., Groenevelt W., Redig F., Orthogonal
 dualities of {M}arkov processes and unitary symmetries, \href{https://doi.org/10.3842/SIGMA.2019.053}{\textit{SIGMA}}
 \textbf{15} (2019), 053, 27~pages, \href{https://arxiv.org/abs/1812.08553}{arXiv:1812.08553}.

\bibitem[8]{carinci2015dualities}
Carinci G., Giardin\`a C., Giberti C., Redig F., Dualities in population
 genetics: a fresh look with new dualities, \href{https://doi.org/10.1016/j.spa.2014.10.009}{\textit{Stochastic Process. Appl.}}
 \textbf{125} (2015), 941--969, \href{https://arxiv.org/abs/1302.3206}{arXiv:1302.3206}.

\bibitem[9]{carinci2021consistent}
Carinci G., Giardin\`a C., Redig F., Consistent particle systems and duality,
 \href{https://doi.org/10.1214/21-ejp684}{\textit{Electron.~J.~Probab.}} \textbf{26} (2021), 125, 31~pages,
 \href{https://arxiv.org/abs/1907.10583}{arXiv:1907.10583}.

\bibitem[12]{DawsonMeasureValuedMarkovProcesses}
Dawson D.A., Measure-valued {M}arkov processes, in \'Ecole d'\'Et\'e de
 {P}robabilit\'es de {S}aint-{F}lour {XXI}---1991, \textit{Lecture Notes in
 Math.}, Vol. 1541, \href{https://doi.org/10.1007/BFb0084190}{Springer}, Berlin, 1993, 1--260.

\bibitem[14]{Etheridge00}
Etheridge A.M., An introduction to superprocesses, \textit{Univ. Lecture Ser.},
 Vol.~20, \href{https://doi.org/10.1090/ulect/020}{American Mathematical Society}, Providence, RI, 2000.

\bibitem[15]{IntertwiningConsistent}
Floreani S., Jansen S., Redig F., Wagner S., Intertwining and duality for
 consistent {M}arkov processes, \href{https://doi.org/10.1214/24-ejp1124}{\textit{Electron.~J.~Probab.}} \textbf{29}
 (2024), 1--34, \href{https://arxiv.org/abs/2112.11885}{arXiv:2112.11885}.

\bibitem[16]{algebraic_Floreani}
Floreani S., Jansen S., Wagner S., Representations of the $\mathfrak{su}(1,1)$
 current algebra, \href{https://arxiv.org/abs/2402.07493}{arXiv:2402.07493}.

\bibitem[17]{floreani2020orthogonal}
Floreani S., Redig F., Sau F., Orthogonal polynomial duality of boundary driven
 particle systems and non-equilibrium correlations, \href{https://doi.org/10.1214/21-aihp1163}{\textit{Ann. Inst. Henri
 Poincar\'e Probab. Stat.}} \textbf{58} (2022), 220--247, \href{https://arxiv.org/abs/2007.08272}{arXiv:2007.08272}.

\bibitem[20]{giardina-redig-vafayi2010}
Giardin\`a C., Redig F., Vafayi K., Correlation inequalities for interacting
 particle systems with duality, \href{https://doi.org/10.1007/s10955-010-0055-0}{\textit{J.~Stat. Phys.}} \textbf{141} (2010),
 242--263.

\bibitem[21]{TransitionMoran}
Karlin S., McGregor J., On a genetics model of {M}oran, \href{https://doi.org/10.1017/S0305004100036513}{\textit{Proc. Cambridge
 Philos. Soc.}} \textbf{58} (1962), 299--311.

\bibitem[23]{kozubowski2008distributional}
Kozubowski T.J., Podg\'orski K., Distributional properties of the negative
 binomial {L}\'evy process, \textit{Probab. Math. Statist.} \textbf{29}
 (2009), 43--71.

\bibitem[27]{LastPenroseLecturesOnThePoissonProcess}
Last G., Penrose M., Lectures on the {P}oisson process, \textit{IMS Textb.},
 Vol.~7, Cambridge University Press, Cambridge, 2018.

\bibitem[31]{MatthesWarmuthMecke}
Matthes K., Warmuth W., Mecke J., Bemerkungen zu einer {A}rbeit: ``{I}ntegral
 and differential characterizations of the {G}ibbs process'', \href{https://doi.org/10.1002/mana.19790880110}{\textit{Math.
 Nachr.}} \textbf{88} (1979), 117--127.

\bibitem[33]{Moran}
Moran P.A.P., Random processes in genetics, \href{https://doi.org/10.1017/s0305004100033193}{\textit{Proc. Cambridge Philos.
 Soc.}} \textbf{54} (1958), 60--71.

\bibitem[35]{Rafler2009}
Rafler M., Gaussian loop- and {P}\'olya processes: a~point process approach,
 Ph.D.~Thesis, {U}niversit\"at {P}otsdam, 2009, available at
 \url{https://publishup.uni-potsdam.de/frontdoor/index/index/docId/3841}.

\bibitem[36]{reedSimonI}
Reed M., Simon B., Methods of modern mathematical physics.~{I}. Functional
 analysis, Academic Press, Sad New York, 1980.

\bibitem[37]{reedSimonII}
Reed M., Simon B., Methods of modern mathematical physics.~{II}. {F}ourier
 analysis, self-adjointness, Academic Press, New York, 1975.

\bibitem[38]{schmudgen2012unbounded}
Schm\"udgen K., Unbounded self-adjoint operators on Hilbert space,
 \textit{Grad. Texts in Math.}, Vol. 265, \href{https://doi.org/10.1007/978-94-007-4753-1}{Springer}, Dordrecht, 2012.

\bibitem[39]{serfozo1990point}
Serfozo R.F., Point processes, in Stochastic Models, \textit{Handbooks Oper.
 Res. Management Sci.}, Vol.~2, \href{https://doi.org/10.1016/S0927-0507(05)80165-3}{North-Holland}, Amsterdam, 1990, 1--93.

\bibitem[40]{shigaDiffusions}
Shiga T., A stochastic equation based on a {P}oisson system for a class of
 measure-valued diffusion processes, \href{https://doi.org/10.1215/kjm/1250520071}{\textit{J.~Math. Kyoto Univ.}} \textbf{30}
 (1990), 245--279.

\end{thebibliography}
\end{document}
